\documentclass[tikz,border=6pt]{standalone}
\usepackage{fontspec}
\usepackage{xeCJK}
\setmainfont{Times New Roman}
\setCJKmainfont{Hiragino Sans GB}
\setCJKsansfont{Microsoft YaHei}
\setCJKmonofont{STSong}
\usepackage{xcolor}

\definecolor{blockblue}{HTML}{4A90D9}
\definecolor{linegray}{HTML}{333333}

\begin{document}
\begin{tikzpicture}[line width=1.2pt]
  \node[
    draw=blockblue,
    rounded corners=4pt,
    inner sep=8pt,
    align=left
  ] at (0,0) {
    \begin{tabular}{p{3.3cm}p{8.4cm}}
      \textbf{条目} & \textbf{结论} \\
      适用范围 & 仅适用于线性定常系统 \\
      模型边界 & 描述输入到输出的关系，不直接呈现内部状态 \\
      单输入单输出 / 多输入多输出 & 一个传递函数只对应一个单输入单输出通道 \\
      系统阶次 & 分母中 $s$ 的最高次幂 \\
      真有理性 & 对物理系统通常有分母阶次 $>$ 分子阶次 \\
      初始状态 & 不能直接表示非零初始条件引起的运动 \\
      特征信息 & $D(s)$ 是特征多项式 \\
    \end{tabular}
  };
\end{tikzpicture}
\end{document}
